\chapter{Action Modeling}
\label{ch:action-modeling}

Everything discussed so far---value compression, preference compression, bias
management---ultimately feeds into one concrete question: \emph{how does the
agent choose an action?}  The answer depends on whether the action space is
discrete or continuous, and whether the policy is deterministic or stochastic.
These choices are not cosmetic; they fundamentally shape the algorithm's
architecture, exploration behavior, and failure modes.


%% ================================================================
\section{Discrete Actions: Enumerate, Then Argmax}
\label{sec:discrete}

When the action space is a finite set $\mathcal{A} = \{a_1, a_2, \ldots, a_k\}$,
the Q-learning approach is natural.  A neural network takes the state as input
and outputs one Q-value per action:
\begin{equation}
  s \;\longrightarrow\; \text{network} \;\longrightarrow\;
  \bigl[Q(s,a_1),\; Q(s,a_2),\; \ldots,\; Q(s,a_k)\bigr].
\end{equation}
Decision-making is trivial: scan the $k$ outputs and pick the largest.
\begin{equation}
  a^{\ast} = \operatorname{argmax}_{a \in \mathcal{A}} Q(s,a).
\end{equation}
DQN is built on exactly this structure.  The convenience of discrete actions is
that argmax is a free operation---$O(k)$ comparisons with no optimization
required.  The limitation is equally clear: the approach does not scale when
$k$ is very large or when actions are continuous.


%% ================================================================
\section{Continuous Actions: The Need for an Actor}
\label{sec:continuous}

In continuous control, actions are real-valued vectors $a \in \mathbb{R}^n$.
Argmax becomes a continuous optimization problem:
\begin{equation}
  a^{\ast} = \operatorname{argmax}_{a \in \mathbb{R}^n} Q(s, a),
\end{equation}
which must be solved at every decision step---generally infeasible in real time.

The solution is to introduce an \textbf{actor network} $\mu_\theta(s)$ that
directly outputs an action, effectively serving as a learned, differentiable
approximation to argmax:
\begin{equation}
  a = \mu_\theta(s), \qquad
  \theta^{\ast} = \operatorname{argmax}_\theta\; Q(s,\, \mu_\theta(s)).
\end{equation}
The actor is trained by backpropagating through the critic's gradient
$\nabla_a Q(s,a)\big|_{a=\mu_\theta(s)}$, pushing its output toward regions
where $Q$ is high.

\begin{keyidea}[The Actor as Differentiable Argmax]
In continuous action spaces, the actor network replaces the intractable
$\operatorname{argmax}_{a \in \mathbb{R}^n} Q(s,a)$ with a learned function
$\mu_\theta(s)$ trained via gradient ascent on $Q$.  DDPG is the canonical
realization of this idea.
\end{keyidea}


%% ================================================================
\section{Policy Representations}
\label{sec:policy-representations}

When the algorithm optimizes a policy directly (as in policy gradient, PPO, or
SAC), the policy must be represented as a parameterized distribution from which
actions can be sampled and whose log-probability can be differentiated.

\subsection{Discrete Policy: Categorical Distribution}

For discrete actions, the standard representation is a categorical distribution
produced by a softmax over network logits:
\begin{equation}
  s \;\longrightarrow\; \text{network} \;\longrightarrow\;
  \ell_1, \ldots, \ell_k
  \;\longrightarrow\;
  \pi(a_i|s) = \frac{\exp(\ell_i)}{\sum_j \exp(\ell_j)}.
\end{equation}
Sampling draws an action from this categorical distribution.  The log-probability
$\log \pi(a|s) = \ell_a - \log\sum_j \exp(\ell_j)$ is smooth and differentiable
with respect to the network parameters.

\subsection{Continuous Policy: Gaussian Distribution}

For continuous actions, the most common choice is a diagonal Gaussian:
\begin{equation}
  \pi(a|s) = \mathcal{N}\!\bigl(a;\; \mu_\theta(s),\; \sigma_\theta(s)^2 I\bigr),
\end{equation}
where the network outputs both the mean $\mu_\theta(s)$ and the standard
deviation $\sigma_\theta(s)$ (or log-std for numerical stability):
\begin{equation}
  s \;\longrightarrow\; \text{network} \;\longrightarrow\;
  (\mu,\; \log\sigma)
  \;\longrightarrow\;
  a = \mu + \sigma \odot \epsilon,
  \quad \epsilon \sim \mathcal{N}(0, I).
\end{equation}
The mean encodes the center of the agent's preference; the standard deviation
encodes the exploration range.  This is the policy representation used by PPO
and SAC in continuous-control tasks.


%% ================================================================
\section{Deterministic vs.\ Stochastic Policies}
\label{sec:det-vs-stoch}

\subsection{Deterministic: Hard Preference}

A deterministic policy $a = \mu_\theta(s)$ outputs a single action per state.
It represents a \emph{hard preference}---full commitment to one choice.

\begin{itemize}[nosep]
  \item \textbf{Strengths}: stable execution, simple inference, no sampling
    variance.
  \item \textbf{Weaknesses}: no intrinsic exploration.  The agent must rely on
    external noise (e.g., Ornstein--Uhlenbeck in DDPG, Gaussian noise in TD3)
    to visit new state--action pairs.
  \item \textbf{Representatives}: DDPG, TD3.
\end{itemize}

\subsection{Stochastic: Soft Preference}

A stochastic policy $a \sim \pi_\theta(a|s)$ outputs a distribution.  It
represents a \emph{soft preference}---a graded ranking of actions rather than a
single commitment.

\begin{itemize}[nosep]
  \item \textbf{Strengths}: built-in exploration, compatible with entropy
    regularization, provides a natural log-probability for policy gradient.
  \item \textbf{Weaknesses}: sampling introduces variance into both training
    and execution.
  \item \textbf{Representatives}: PPO, SAC, REINFORCE.
\end{itemize}

Table~\ref{tab:det-stoch} summarizes the comparison.

\begin{table}[ht]
\centering
\caption{Deterministic vs.\ stochastic policies.}
\label{tab:det-stoch}
\begin{tabular}{@{}lll@{}}
\toprule
& \textbf{Deterministic} & \textbf{Stochastic} \\
\midrule
Output       & Single action $\mu(s)$ & Distribution $\pi(a|s)$ \\
Preference   & Hard (point estimate)  & Soft (distribution) \\
Exploration  & External noise required & Intrinsic \\
Gradient signal & $\nabla_a Q$ through actor & $\nabla\log\pi \cdot A$ \\
Typical algorithms & DDPG, TD3 & PPO, SAC, REINFORCE \\
\bottomrule
\end{tabular}
\end{table}


%% ================================================================
\section{Why RL Moved Toward Stochastic Policies}
\label{sec:why-stochastic}

In classical control, deterministic policies are the norm: given a state, output
the optimal action.  Why did RL increasingly favor stochastic policies?

The answer lies in RL's unique structure: the policy serves \textbf{two roles
simultaneously}.

\begin{enumerate}
  \item \textbf{Decision-maker}: choosing what to do right now.
  \item \textbf{Data sampler}: determining what the agent will observe in the
    future.
\end{enumerate}

A deterministic policy is a competent decision-maker but a poor data sampler.
It concentrates all future data around a single trajectory, starving the
learning algorithm of the diversity it needs to improve.

A stochastic policy balances both roles.  It makes good (on average) decisions
\emph{and} maintains the data coverage needed for continued learning.

\begin{keyidea}[The Dual Role of Policy]
In RL, the policy is both the decision rule and the data-generation mechanism.
A stochastic policy is a soft preference representation that naturally balances
exploitation (good decisions) with exploration (diverse data).  This dual role
is the fundamental reason RL moved from deterministic to stochastic policies.
\end{keyidea}

From the compression perspective of Chapter~\ref{ch:compression}, a stochastic
policy is a richer, more complete compression of preferences.  Preferences are
not always a single point; they can---and often should---be a distribution that
reflects the agent's uncertainty about which action is truly best.
\begin{align}
  \text{Deterministic policy} &= \text{hard preference (point compression)}, \\
  \text{Stochastic policy}    &= \text{soft preference (distributional compression)}.
\end{align}
The shift from deterministic to stochastic is, at its core, a shift from lossy
point compression to richer distributional compression of behavioral preferences.
